\documentclass[11pt]{article}
\usepackage{iftex}
\ifPDFTeX
  \usepackage[utf8]{inputenc}\usepackage[T5]{fontenc}
\else
  \usepackage{fontspec}
\fi
\usepackage[a4paper,top=2.95cm,headheight=1.5cm,headsep=0.45cm,bottom=3.1cm,left=1.9cm,right=1.9cm]{geometry}
\usepackage{amsmath,amssymb}
\usepackage{tikz}\usetikzlibrary{calc,shadows,shapes.geometric}
\usepackage[most]{tcolorbox}
\usepackage{enumitem}
\usepackage{array}
\usepackage{tabularx}
\usepackage{colortbl}
\usepackage{comment}
\usepackage{needspace}
\usepackage{environ}
\usepackage{fancyhdr}
\usepackage{fontawesome5}
\usepackage{graphicx}
\usepackage{microtype}
% PDF outline (bookmarks) so each chapter shows in the reader's outline — does not change the page layout
\usepackage[hidelinks,unicode,bookmarksnumbered=false,bookmarksopen=true,bookmarksopenlevel=0]{hyperref}
\usepackage{bookmark}
\usetikzlibrary{decorations.pathreplacing,arrows.meta,angles,quotes,patterns,positioning,through,intersections}

% ---------------- palette: navy family + subtle light tint; orange ONLY in the logo ----------------
% ----- EXACT palette from the user's original main (36).tex (vivid, not pale) -----
\definecolor{tmBlue}{RGB}{66,133,244}   % problem / header / section
\definecolor{tmGreen}{RGB}{52,168,83}   % solution
\definecolor{tmOrange}{RGB}{255,152,0}  % example
\colorlet{Navy}{tmBlue}                 % header rule / section heading / badge / divider
\definecolor{NavyD}{HTML}{222222}       % solution body text (dark, readable)
\colorlet{Gold}{Navy}
\colorlet{ProbBlue}{tmBlue}             % problem boxes
\colorlet{SolGreen}{tmGreen}            % solution boxes
\definecolor{AnsRed}{HTML}{E02424}      % vivid red for answers / totals
\definecolor{CardGreen}{HTML}{34A853}   % digit-card green
\definecolor{StarGold}{HTML}{F2B134}    % difficulty star (gold)
\definecolor{KidOrange}{HTML}{DE8A2C}   % logo accent only
\definecolor{Ink}{HTML}{1F2733}       % body text
\definecolor{Slate}{HTML}{5B6A82}     % muted labels
\definecolor{Line}{HTML}{C6D0DF}      % light hairlines / student boxes
\definecolor{Tint}{HTML}{EDF1F7}      % pale panel fill (worked solution / answer)
% aliases so older names keep working
\colorlet{TextDark}{Ink}\colorlet{TextGray}{Slate}\colorlet{Soft}{Slate}
\colorlet{CardLine}{Line}\colorlet{Teal}{Navy}\colorlet{ProbB}{Navy}\colorlet{SolG}{Navy}
% bright, kid-friendly colours for bar-model (sơ đồ đoạn thẳng) blocks
\definecolor{tmPurple}{RGB}{124,77,226}
\definecolor{tmTeal}{RGB}{0,150,136}
\definecolor{tmPink}{RGB}{233,30,99}
\AtBeginDocument{\color{Ink}}
\raggedbottom  % natural flow, even spacing (no stretched gaps)
\setlength{\parindent}{0pt}\setlength{\parskip}{0.18em}
\renewcommand{\arraystretch}{1.35}
\providecommand{\AnsWord}{Đáp số}

% ---------------- brand logo (original mark) ----------------
\newcommand{\logomark}[1]{\begin{tikzpicture}[baseline=-0.55ex,scale=#1]
  \draw[KidOrange,line width=1.5pt] (-0.27,0) .. controls (-0.27,0.23) and (0,0.23) .. (0,0)
        .. controls (0,-0.23) and (-0.27,-0.23) .. (-0.27,0);
  \draw[Navy,line width=1.5pt] (0.27,0) .. controls (0.27,0.23) and (0,0.23) .. (0,0)
        .. controls (0,-0.23) and (0.27,-0.23) .. (0.27,0);
\end{tikzpicture}}
% real brand logo image (gold infinity + "think algo" wordmark)
\newcommand{\logohead}{\IfFileExists{logo.png}{\includegraphics[height=1.35cm]{logo.png}}%
  {\raisebox{0.35cm}{\logomark{1.5}\hspace{0.5em}{\Large\bfseries\color{KidOrange}think\,\color{Navy}algo}}}}

% ---------------- running header (logo + Name/Date + topic) on EVERY page, single navy rule ----------------
\pagestyle{fancy}\fancyhf{}
\renewcommand{\headrulewidth}{1.4pt}\renewcommand{\footrulewidth}{0pt}
\renewcommand{\headrule}{{\color{Navy}\hrule height 1.3pt}}
\fancyhead[L]{\raisebox{-0.95cm}{\logohead}}
\fancyhead[R]{\setlength{\tabcolsep}{0pt}%
  \begin{tabular}{@{}r@{}}
    {\small\color{Ink}{\color{tmBlue}\faUserEdit}~\NameWord:~\rule{3.0cm}{0.5pt}\quad{\color{tmBlue}\faCalendarDay}~\DateWord:~\rule{1.9cm}{0.5pt}}\\[5pt]
    {\large\bfseries\color{Navy}\CurTopic}
  \end{tabular}}

% ---------------- footer (rounded page border + Topic/Grade + thinkalgo.com + page + tagline) ----------------
% NOTE: a remember-picture overlay must be placed DIRECTLY in \fancyfoot, NOT inside \def/\newcommand.
\fancyfoot[C]{\begin{tikzpicture}[remember picture, overlay]
  \draw[Navy!40, line width=1.4pt, rounded corners=12pt]
     ($(current page.north west)+(0.8cm,-0.8cm)$) rectangle ($(current page.south east)+(-0.8cm,0.45cm)$);
  \coordinate (FL) at ($(current page.south west)+(1.2cm, 2.18cm)$);
  \coordinate (FR) at ($(current page.south east)+(-1.2cm, 2.18cm)$);
  \draw[Navy!55, line width=1pt] (FL) -- (FR);
  \node[fill=white, draw=Navy, line width=1.2pt, rounded corners=8pt, font=\bfseries\small,
        text=Navy, inner xsep=11pt, inner ysep=4pt] at ($(current page.south)+(0,2.18cm)$) {{\color{tmBlue}\faGlobe}\ \ thinkalgo.com};
  \node[text=Slate, font=\footnotesize, anchor=north west, align=left, inner sep=0pt, text width=8.2cm]
        at ($(FL)+(0.1cm,-0.30cm)$) {{\color{tmBlue}\faTag}\ \TopicWord: \textbf{\color{Ink}\CurTopic}\\[1.5pt]{\color{tmBlue}\faGraduationCap}\ \DocGrade};
  \node[text=Slate, font=\footnotesize] at ($(current page.south)+(0,1.20cm)$) {\thepage};
  \node[text=Slate, font=\scriptsize, anchor=north east, align=right, inner sep=0pt, text width=7.4cm]
        at ($(FR)+(-0.1cm,-0.34cm)$) {{\normalsize\bfseries\color{Navy}ThinkAlgo}\ — From Math Thinking to Code};
\end{tikzpicture}}

% ---------------- page-1 lead-in (topic is already in the header; keep a slim spacer) ----------------
\newcommand{\titlehead}{\vspace*{-0.2cm}}
\providecommand{\CurTopic}{}

% chapter banner + sets the running-header topic. #1 = number, #2 = title
\newcommand{\chapterhead}[2]{\renewcommand{\CurTopic}{#2}%
  \setcounter{seccnt}{0}% each chapter is independent: section badges restart at 1
  \phantomsection\pdfbookmark[0]{#1. #2}{chap#1}\label{chap:#1}% add this chapter to the PDF outline + TOC ref
  \par\vspace{0.1cm}\noindent\begin{tikzpicture}[baseline=-0.5ex]
    \node[circle,fill=tmBlue,text=white,font=\Large\bfseries,minimum size=0.95cm,inner sep=0pt](cn){#1};
    \node[anchor=west,font=\LARGE\bfseries,text=tmBlue] at ([xshift=12pt]cn.east){#2};
  \end{tikzpicture}\par\vspace{4pt}\noindent\textcolor{tmBlue}{\rule{\linewidth}{2pt}}\par\vspace{0.5cm}}

% ---------------- section heading (outline navy number badge, no fill) ----------------
\newcounter{seccnt}
\newcommand{\shead}[1]{\par\vspace{0.7em}\refstepcounter{seccnt}%
  \noindent\begin{tikzpicture}[baseline=-0.6ex]
    \node[circle,draw=Gold,line width=1.2pt,text=Gold,font=\bfseries,minimum size=1.05em,inner sep=1pt]{\theseccnt};
  \end{tikzpicture}\hspace{0.4em}{\large\bfseries\color{Navy} #1}\par\vspace{0.25em}}

% ============ boxes EXACTLY per the original main (36).tex: solid-blue PROBLEM, solid-orange EXAMPLE,
%              outline-green SOLUTION, all with drop shadow + rounded corners ============
% ---------------- knowledge box: white, tmBlue border + left accent + drop shadow ----------------
\newtcolorbox{knowbox}{enhanced,breakable,colback=white,colframe=tmBlue!70,boxrule=0.9pt,arc=5pt,
  left=12pt,right=11pt,top=6pt,bottom=6pt,
  borderline west={3.5pt}{0pt}{tmBlue}}

% ---------------- EXAMPLE (Ví dụ): SOLID ORANGE tab, white text ----------------
\newtcolorbox{exbox}[1]{enhanced,breakable,colback=white,colframe=tmOrange,boxrule=1pt,arc=5pt,
  left=12pt,right=11pt,top=14pt,bottom=7pt,
  attach boxed title to top left={xshift=10pt,yshift=-\tcboxedtitleheight/2},
  boxed title style={colback=tmOrange,colframe=tmOrange,arc=4pt,boxrule=0pt},
  title={#1},fonttitle=\bfseries\color{white}}
\newcommand{\guide}{\par\vspace{2pt}{\bfseries\color{tmBlue}\GuideWord}\par\nopagebreak}

% ---------------- PROBLEM (Bài): SOLID BLUE tab, white text ----------------
\newtcolorbox{probboxT}[2]{enhanced,breakable,colback=white,colframe=tmBlue,boxrule=1pt,arc=5pt,
  left=12pt,right=10pt,top=14pt,bottom=7pt,
  attach boxed title to top left={xshift=10pt,yshift=-\tcboxedtitleheight/2},
  boxed title style={colback=tmBlue,colframe=tmBlue,arc=4pt,boxrule=0pt},
  fonttitle=\bfseries\color{white},title={\ProbWord\ #1\ifx\relax#2\relax\else#2\fi}}
% unbreakable copy used inside the student block
\newtcolorbox{probboxS}[2]{enhanced,colback=white,colframe=tmBlue,boxrule=1pt,arc=5pt,
  left=12pt,right=10pt,top=14pt,bottom=7pt,
  attach boxed title to top left={xshift=10pt,yshift=-\tcboxedtitleheight/2},
  boxed title style={colback=tmBlue,colframe=tmBlue,arc=4pt,boxrule=0pt},
  fonttitle=\bfseries\color{white},title={\ProbWord\ #1\ifx\relax#2\relax\else#2\fi}}
% Teacher: the problem box is printed where it stands.
% Student: the problem is stored and printed by \probsol together with its blank "work" box inside one
%          unbreakable block, so a problem is never separated from its writing space by a page break.
\ifsol
  \newenvironment{probbox}[2]{\begin{probboxT}{#1}{#2}}{\end{probboxT}}
\else
  \NewEnviron{probbox}[2]{\global\let\StuProbBody\BODY\gdef\StuProbArgs{{#1}{#2}}}
  \newcommand{\StuProb}[2]{\begin{probboxS}{#1}{#2}\StuProbBody\end{probboxS}}
\fi

% ---------------- SOLUTION (Lời giải): OUTLINE GREEN tab (white fill, green border + text) ----------------
% Teacher: a green tcolorbox. Student: the body is dropped entirely (comment pkg — robust with tables/\\).
\ifsol
  \newtcolorbox{solbox}{enhanced,breakable,colback=white,colframe=tmGreen,boxrule=1pt,arc=5pt,
    left=12pt,right=10pt,top=14pt,bottom=7pt,
    attach boxed title to top left={xshift=10pt,yshift=-\tcboxedtitleheight/2},
    boxed title style={colback=white,colframe=tmGreen,boxrule=1pt,arc=4pt},
    fonttitle=\bfseries\color{tmGreen!88!black},title={\SolTitle},before upper={\color{NavyD}}}
\else
  \excludecomment{solbox}
\fi

% ---------------- answer highlight: label green, value RED ----------------
\newcommand{\ans}[1]{\par\smallskip{\bfseries\color{tmGreen!88!black}\AnsWord:}\ {\bfseries\color{AnsRed}#1}\par}

% ---------------- digit cards (exact original style: fill=c!10, border=c!60!black, drop shadow) ----------------
% \digitcards{0/gray,1/cyan,2/green,3/cyan} : value/colour pairs
\newcommand{\digitcards}[1]{\begin{center}\begin{tikzpicture}[scale=0.7,transform shape,font=\sffamily]
  \def\dx{0}\foreach \v/\c in {#1}{
    \node[draw=\c!60!black,line width=1.5pt,fill=\c!10,rounded corners=4pt,
          minimum width=1cm,minimum height=1.4cm] at (\dx,0){\Huge\bfseries \v};
    \pgfmathsetmacro\dx{\dx+1.5}\xdef\dx{\dx}}
\end{tikzpicture}\end{center}}

% ---------------- student working space (white, light border) ----------------
\newcommand{\worksp}[1]{\ifsol\else\begin{tcolorbox}[colback=white,colframe=Line,boxrule=0.8pt,
  arc=5pt,height=\dimexpr#1+0.05cm\relax,left=8pt,right=8pt,top=4pt,bottom=4pt]{\footnotesize\color{Slate}\itshape\WorkWord}\end{tcolorbox}\fi}

% ---------------- \probsol{solution} : Teacher shows the green solution box; Student shows a BLANK
%   "Bài làm" box whose height matches that solution's height (writing room = work needed), so both
%   versions have identical per-problem height and paginate the same. ----------------
\newsavebox{\solmeasure}\newlength{\solht}
\newcommand{\probsol}[1]{%
  \begin{lrbox}{\solmeasure}\begin{minipage}{\dimexpr\linewidth-24pt\relax}#1\end{minipage}\end{lrbox}%
  \ifsol
    \begin{solbox}#1\end{solbox}%
  \else
    \setlength{\solht}{\dimexpr\ht\solmeasure+\dp\solmeasure+0.85cm\relax}%
    \ifdim\solht<1.35cm\setlength{\solht}{1.35cm}\fi
    \ifdim\solht>\StudentMaxWork\setlength{\solht}{\StudentMaxWork}\fi
    \par\noindent\begin{minipage}{\linewidth}%
    \expandafter\StuProb\StuProbArgs
    \begin{tcolorbox}[colback=white,colframe=Line,boxrule=0.8pt,arc=5pt,
      left=8pt,right=8pt,top=4pt,bottom=4pt,height=\solht]{\footnotesize\color{Slate}\itshape\WorkWord}\end{tcolorbox}%
    \end{minipage}\par\medskip
  \fi}

% ---------------- practice divider: blue left accent bar, no fill ----------------
\newcommand{\dividerbar}[1]{\par\vspace{0.6em}\noindent
  \begin{tcolorbox}[enhanced,frame hidden,interior hidden,borderline west={5pt}{0pt}{ProbBlue},
    colback=white,left=12pt,right=12pt,top=5pt,bottom=5pt,arc=4pt]
  {\large\bfseries\color{ProbBlue} #1}\end{tcolorbox}\vspace{0.4em}}

% place-value choice-count diagram (Cambridge style): #1 = "label/count/note/color" list, #2 = result
% notes must NOT contain commas or slashes (they separate \foreach fields)
\newcommand{\countboxes}[2]{\begin{center}\begin{tikzpicture}[scale=0.62,transform shape,font=\sffamily]
  \def\xx{0}%
  \foreach \lab/\cnt/\note/\col [count=\i] in {#1}{%
    \ifnum\i>1\relax \node[text=gray,font=\Huge] at (\xx-1.8,0){$\times$};\fi
    \node[draw=\col!65,fill=\col!8,rounded corners=3pt,line width=1pt,minimum width=2.7cm,minimum height=1.7cm] at (\xx,0){};
    \node[text=\col!85,font=\bfseries] at (\xx,0.58){\lab};
    \node[text=\col!88,font=\Huge\bfseries] at (\xx,-0.1){\cnt};
    \node[text=gray,font=\scriptsize] at (\xx,-1.05){\note};
    \pgfmathsetmacro\xx{\xx+3.6}\xdef\xx{\xx}}%
  \node[text=gray,font=\Huge] at (\xx-1.8,0){$=$};
  \node[text=red,font=\Huge\bfseries] at (\xx,0){#2};
\end{tikzpicture}\end{center}}

% ======== bar model (sơ đồ đoạn thẳng): lightly-rounded blocks, darker border, tiny gaps, bright colours ========
% one block: \blk{x1}{x2}{yc}{height}{colour}{label}  (a small gap is left on each side)
\newcommand{\blk}[6]{%
  \draw[rounded corners=1.3pt,fill=#5!16,draw=#5!80!black,line width=1pt]
    (#1+0.045,#3-#4/2) rectangle (#2-0.045,#3+#4/2);
  \node[font=\small\bfseries,text=#5!52!black] at ({(#1+#2)/2},#3){#6};}
% measurement with a double-headed ARROW (not a brace) + label above the span: \brcup{x1}{x2}{y}{label}
\newcommand{\brcup}[4]{%
  \draw[<->,>=stealth,tmBlue!72!black,line width=0.7pt] (#1+0.05,#3)--(#2-0.05,#3);
  \node[above=-1.5pt,font=\footnotesize,text=tmBlue!58!black] at ({(#1+#2)/2},#3){#4};}
% measurement with a double-headed ARROW + label below the span: \brcdn{x1}{x2}{y}{label}
\newcommand{\brcdn}[4]{%
  \draw[<->,>=stealth,tmBlue!72!black,line width=0.7pt] (#1+0.05,#3)--(#2-0.05,#3);
  \node[below=-1.5pt,font=\footnotesize,text=tmBlue!58!black] at ({(#1+#2)/2},#3){#4};}
% row label at the left of a bar: \rowlab{x}{y}{text}
\newcommand{\rowlab}[3]{\node[anchor=east,font=\small\bfseries,text=Ink] at (#1,#2){#3};}
% vertical right-pointing CURLY BRACE spanning two stacked bars: \vbrace{x}{ytop}{ybot}{label}
\newcommand{\vbrace}[4]{%
  \pgfmathsetmacro\ym{(#2+#3)/2}%
  \draw[tmBlue!72!black,line width=0.9pt,line cap=round]
    (#1,#2) .. controls (#1+0.18,#2) and (#1+0.18,#2-0.18) .. (#1+0.18,#2-0.22)
    -- (#1+0.18,{\ym+0.22})
    .. controls (#1+0.18,\ym) and (#1+0.30,{\ym+0.08}) .. (#1+0.42,\ym)
    .. controls (#1+0.30,{\ym-0.08}) and (#1+0.18,\ym) .. (#1+0.18,{\ym-0.22})
    -- (#1+0.18,#3+0.22)
    .. controls (#1+0.18,#3+0.18) and (#1+0.18,#3) .. (#1,#3);
  \node[right=4pt,font=\footnotesize,text=tmBlue!58!black] at (#1+0.42,\ym){#4};}
% blank 3x3 grid for a problem statement (student fills it in)
\newcommand{\blankgrid}{\begin{tikzpicture}[font=\sffamily]
  \foreach \r in {0,1,2}{\foreach \c in {0,1,2}{
    \draw[rounded corners=1.3pt,fill=Tint,draw=Slate!55,line width=0.9pt]
      (\c+0.05,\r+0.05) rectangle (\c+0.95,\r+0.95);}}
\end{tikzpicture}}
% wrapper: LEFT-aligned (not centred), a little air above/below
\newenvironment{barfig}{\par\vspace{3pt}\noindent\kern0.35cm\begin{tikzpicture}[font=\sffamily,>=stealth]}%
  {\end{tikzpicture}\par\vspace{3pt}}
% wrapper: CENTRED (for illustrative diagrams in theory / examples)
\newenvironment{barfigC}{\par\vspace{3pt}\begin{center}\begin{tikzpicture}[font=\sffamily,>=stealth]}%
  {\end{tikzpicture}\end{center}\vspace{2pt}}

% ---------------- difficulty stars (nicely drawn) : \stars{1}=hard, \stars{2}=very hard ----------------
\newcommand{\diffstar}{\tikz[baseline=-0.55ex]{\node[star,star points=5,star point ratio=0.46,
  draw=StarGold!72!black,fill=StarGold,line width=0.3pt,inner sep=1.75pt]{};}}
\newcommand{\stars}[1]{\ifcase#1\or\kern2pt\diffstar\or\kern2pt\diffstar\kern1.5pt\diffstar\fi}

\newcommand{\nd}[1]{\overline{#1}}

% number line: #1 last index, #2 ticks "i/label", #3 open (excluded) indices, #4 filled (solution) indices
\newcommand{\numlinex}[4]{\begin{center}\begin{tikzpicture}[x=1.55cm,>=stealth]
  \draw[->,thick,Navy](-0.45,0)--(#1+0.45,0);
  \foreach \i/\lab in {#2}{\draw[Navy,line width=0.8pt](\i,0.08)--(\i,-0.08);
     \node[below,font=\small,text=Ink]at(\i,-0.1){\lab};}
  \foreach \i in {#3}{\draw[Navy,fill=white,line width=1pt](\i,0)circle(2.3pt);}
  \foreach \i in {#4}{\fill[KidOrange](\i,0)circle(2.7pt);}
\end{tikzpicture}\end{center}}


% ======================================================================
% Grade 7 additions (shared by all four editions)
% ======================================================================
% student edition: blank "work" box height = min(solution height, this cap)
\newlength{\StudentMaxWork}\setlength{\StudentMaxWork}{4.5cm}

% ThinkAlgo tip box (label set per language by \TipWord)
\newtcolorbox{tipbox}{enhanced,breakable,colback=StarGold!8,colframe=StarGold!75!black,boxrule=1pt,arc=4pt,
  left=11pt,right=10pt,top=6pt,bottom=6pt,
  before upper={\textbf{\color{StarGold!65!black}\faLightbulb\ \TipWord}\ }}

% "Recall" box: re-states a piece of theory inside a solution / example
\newtcolorbox{recallbox}{enhanced,breakable,colback=tmBlue!5,colframe=tmBlue!45,boxrule=0.6pt,arc=3pt,
  left=8pt,right=8pt,top=4pt,bottom=4pt,
  before upper={\textbf{\color{tmBlue!80!black}\faBookOpen\ \RecallWord}\ }}

% "Why is it true? / Proof" box
\newtcolorbox{proofbox}{enhanced,breakable,colback=tmPurple!4,colframe=tmPurple!45,boxrule=0.6pt,arc=3pt,
  left=8pt,right=8pt,top=4pt,bottom=4pt,
  before upper={\textbf{\color{tmPurple!85!black}\faSearch\ \ProofWord}\ },
  after upper={\unskip\nobreak\hfill\penalty50\hskip1em\null\nobreak\hfill$\square$}}

% "Common mistake" box
\newtcolorbox{warnbox}{enhanced,breakable,colback=AnsRed!4,colframe=AnsRed!45,boxrule=0.6pt,arc=3pt,
  left=8pt,right=8pt,top=4pt,bottom=4pt,
  before upper={\textbf{\color{AnsRed!85!black}\faExclamationTriangle\ \WarnWord}\ }}

% "Note for parents" box (short hint on how to explain the idea at home)
\newtcolorbox{parentbox}{enhanced,breakable,colback=tmTeal!5,colframe=tmTeal!50,boxrule=0.6pt,arc=3pt,
  left=8pt,right=8pt,top=4pt,bottom=4pt,
  before upper={\textbf{\color{tmTeal!80!black}\faUsers\ \ParentWord}\ }}

% numbered step inside a solution:  \step{1}{Short title}  -> "Bước 1. Short title"
\newcommand{\step}[2]{\par\smallskip\noindent{\bfseries\color{tmBlue!85!black}\StepWord~#1.}\ {\bfseries #2}\par\nopagebreak}
% verification line:  \verify{...}
\newcommand{\verify}[1]{\par\smallskip\noindent{\bfseries\color{tmGreen!80!black}\faCheckCircle\ \CheckWord}\ #1\par}

% multiple-choice options
% \choices{..}{..}{..}{..}{..}      : five short options on one line
% \choicesII{..}{..}{..}{..}{..}    : two/three per line (medium length)
% \choicesL{..}{..}{..}{..}{..}     : one option per line (long text)
\newcommand{\choices}[5]{\par\smallskip\noindent
  \begin{tabularx}{\linewidth}{@{}XXXXX@{}}(A)\ #1&(B)\ #2&(C)\ #3&(D)\ #4&(E)\ #5\end{tabularx}\par}
\newcommand{\choicesII}[5]{\par\smallskip\noindent
  \begin{tabularx}{\linewidth}{@{}XXX@{}}(A)\ #1&(B)\ #2&(C)\ #3\\(D)\ #4&(E)\ #5&\end{tabularx}\par}
\newcommand{\choicesL}[5]{\par\smallskip\noindent
  \begin{tabular}{@{}l@{}}(A)\ #1\\(B)\ #2\\(C)\ #3\\(D)\ #4\\(E)\ #5\end{tabular}\par}

% centred figure inside a problem / solution
\newenvironment{figC}{\par\vspace{2pt}\begin{center}\begin{tikzpicture}[font=\sffamily\small,>=stealth,line join=round]}%
  {\end{tikzpicture}\end{center}\vspace{1pt}}
